\aufzaehlungdrei{Matriks Jacobi tersebut adalah

\mathdisp {\begin{pmatrix}  0   &   1  &  0  \\  -1 &  0  &  0 \\ 0  &  0 &  0  \end{pmatrix}} { . }

}{Sebagai medan normal, kita ambil $\begin{pmatrix}  x  \\ y \\ z   \end{pmatrix}$. Di

\mavergleichskette
{\vergleichskette
{ P
}
{ = }{ \begin{pmatrix}  0  \\ 1 \\  0   \end{pmatrix}
}
{  }{
}
{    }{
}
{   }{
}
}
{}{}{,}
vektor-vektor
\mathkor {} {\begin{pmatrix}  1  \\ 0 \\  0   \end{pmatrix}} {dan} {\begin{pmatrix}  0  \\ 0\\  1   \end{pmatrix}} {}
membentuk suatu basis ruang singgung. Kita mempunyai

\mavergleichskettedisp
{\vergleichskette
{ \begin{pmatrix}  0   &   1  &  0  \\  -1 &  0  &  0 \\ 0  &  0 &  0  \end{pmatrix} \begin{pmatrix}  1  \\ 0 \\  0   \end{pmatrix}
}
{ =} { \begin{pmatrix}  0  \\ -1\\  0   \end{pmatrix}
}
{ } {
}
{ } {
}
{ } {
}
}
{}{}{}
dan

\mavergleichskettedisp
{\vergleichskette
{  \begin{pmatrix}  0   &   1  &  0  \\  -1 &  0  &  0 \\ 0  &  0 &  0  \end{pmatrix} \begin{pmatrix}  0  \\ 0\\  1   \end{pmatrix}
}
{ =} { \begin{pmatrix}  0  \\ 0\\  0   \end{pmatrix}
}
{ } {
}
{ } {
}
{ } {
}
}
{}{}{.}
Di sini,
\mathl{\begin{pmatrix}  0  \\ -1\\  0   \end{pmatrix}}{} merupakan kelipatan vektor normal, sehingga proyeksi ortogonalnya ke ruang singgung juga merupakan vektor nol. Dengan demikian, pemetaan
\mathl{v \mapsto \nabla_vF}{} adalah pemetaan nol.
}{Di

\mavergleichskette
{\vergleichskette
{ Q
}
{ = }{ \begin{pmatrix}  0  \\ 0\\  1   \end{pmatrix}
}
{  }{
}
{    }{
}
{   }{
}
}
{}{}{,}
vektor-vektor
\mathkor {} {\begin{pmatrix}  1  \\ 0 \\  0   \end{pmatrix}} {dan} {\begin{pmatrix}  0  \\ 1 \\  0   \end{pmatrix}} {}
membentuk suatu basis ruang singgung. Kita mempunyai

\mavergleichskettedisp
{\vergleichskette
{ \begin{pmatrix}  0   &   1  &  0  \\  -1 &  0  &  0 \\ 0  &  0 &  0  \end{pmatrix} \begin{pmatrix}  1  \\ 0 \\  0   \end{pmatrix}
}
{ =} { \begin{pmatrix}  0  \\ -1\\  0   \end{pmatrix}
}
{ } {
}
{ } {
}
{ } {
}
}
{}{}{}
dan

\mavergleichskettedisp
{\vergleichskette
{ \begin{pmatrix}  0   &   1  &  0  \\  -1 &  0  &  0 \\ 0  &  0 &  0  \end{pmatrix} \begin{pmatrix}  0  \\ 1 \\  0   \end{pmatrix}
}
{ =} { \begin{pmatrix}  1  \\ 0 \\  0   \end{pmatrix}
}
{ } {
}
{ } {
}
{ } {
}
}
{}{}{.}
Kedua vektor citra ini sudah tangensial pada $Y$ di $Q$. Suatu matriks yang mendeskripsikan
\mathl{v \mapsto \nabla_vF}{} adalah

\mathdisp {\begin{pmatrix}  0   &   1  \\  -1  &  0 \end{pmatrix}} { . }

\emph{Catatan edisi: Tanda pada kedua kolom diperbaiki agar matriks ini sesuai dengan citra kedua vektor basis yang dihitung tepat di atasnya.}

}
